\section{Conclusion}
\label{sec:conclusion}

This paper establishes a disturbance-aware benchmark reading of dynamic frequency estimation for low-inertia IBR grids. The central result is clear: no single estimator dominates once the benchmark spans clean standards-style tests, IBR harmonics, noisy ringdown, and composite multi-event stress. SOGI-PLL, EKF, RA-EKF, SOGI-FLL, UKF, PLL, ESPRIT, and Koopman (RK-DPMU) each lead at least one disturbance family on RMSE, while the trip-risk leaders often differ from the RMSE leaders. The benchmark therefore resolves not to one global winner, but to a structured map of conditional strengths and weaknesses.

Three conclusions follow. First, a latency-robustness trade-off remains structural rather than accidental: methods that suppress spectral pollution through longer memory or heavier filtering often pay for it through slower transient recovery. Second, compliance-style scenarios are necessary but insufficient. The large phase-jump and composite IBR cases expose failure modes that are largely invisible in cleaner standard families. Third, deployment recommendations must remain objective-aware. If the design priority is low-cost ramp tracking, loop methods remain competitive. If the estimator must remain broadly trustworthy under mixed disturbances, the model-based filters are the strongest default recommendation in this study. If the priority is spectral precision or post-event analysis, some window-based methods remain attractive despite their higher delay.

The methodological consequence is equally important. Benchmark conclusions should report the disturbance family, metric contract, and deployment objective explicitly; otherwise, a single ranking can hide the very behavior the benchmark was designed to expose. That disturbance-aware, objective-aware reading is the main contribution of the study.
